\section{The Preview State Machine}
\label{sec:state}

\subsection{States and events}

For a selected entry the preview pane is in one of the states
\[
  \mathsf{idle},\ \mathsf{loading},\ \mathsf{folder},\ \mathsf{image},\
  \mathsf{quicklook},\ \mathsf{pdf},\ \mathsf{archive},\ \mathsf{text},\ \mathsf{binary},\
  \mathsf{empty},\ \mathsf{dataless},\ \mathsf{error}.
\]
The name of the entry gives a \emph{hint} $h\in\{\text{quicklook, image, pdf,
archive, none}\}$; it is only a hint because the server decides (from the bytes,
or for Quick Look from what Quick Look can draw).  Selecting
an entry moves to $\mathsf{loading}$, after a short delay that keeps holding an
arrow key from firing a request per row, and then:
\begin{itemize}
  \item with hint quicklook (checked first, because an iWork package lists as a
        folder), and the entry is not a dangling link, the pane shows
        $\mathsf{quicklook}$ and, if the picture fails to load, goes to
        $\mathsf{folder}$ for a package or falls back to the head request for a
        file;
  \item a folder goes to $\mathsf{folder}$; a dangling link to $\mathsf{error}$;
  \item with hint image, the pane shows $\mathsf{image}$ and, if the server
        refuses the bytes (415), falls back to the head request;
  \item with hint pdf, a one-byte range request asks the server; a
        \code{200}/\code{206} of type \code{application/pdf} gives
        $\mathsf{pdf}$, a \code{409} gives $\mathsf{dataless}$, anything else
        falls back to the head request;
  \item with hint archive, the listing request gives $\mathsf{archive}$, a
        \code{409} gives $\mathsf{dataless}$, anything else falls back to the head
        request;
  \item otherwise the head request gives one of $\mathsf{text}$,
        $\mathsf{binary}$, $\mathsf{empty}$, $\mathsf{dataless}$, or, on failure,
        $\mathsf{error}$.
\end{itemize}

\subsection{Laws}

\begin{law}[Termination]
\label{law:terminate}
From $\mathsf{loading}$ every run reaches a terminal state
(folder, image, quicklook, pdf, archive, text, binary, empty, dataless, or error)
after at most two requests: one hinted request and at most one fall back to the
head request, which never falls back again.
\end{law}

\begin{law}[Freshness]
\label{law:fresh}
The content shown always belongs to the entry that is currently selected: every
request started for a selection, including the fall back after an image or a Quick
Look picture fails to load, is cancelled when the selection changes, and a cancelled request never
changes the state.
\end{law}

\begin{law}[Frame only after confirmation]
\label{law:frame-confirmed}
An inline frame is created for a PDF only in state $\mathsf{pdf}$, which is
entered only after the server has answered \code{application/pdf} for that
path; a Markdown frame is created only for a text head whose name is a Markdown
name, and receives only sanitized HTML.
\Tests{e2e: ``shows a real PDF in the built-in viewer frame, and treats a fake .pdf as text''; Markdown tests}
\end{law}

\begin{law}[Nothing is shown for a denied entry]
\label{law:preview-denied}
Every request of the state machine goes through the endpoints of
Section~\ref{sec:access}; for a denied entry each answers as for a missing one
(Law~\ref{law:indist}), so the pane shows only ``Not available''.
\Tests{e2e: never lists secret files; answers 404}
\end{law}

\begin{boundary}[Termination and freshness are argued, not enumerated]
Laws~\ref{law:terminate} and~\ref{law:fresh} are properties of the component's
control flow.  Freshness was violated by one path (Finding~\ref{find:stale}) and
is now guaranteed by keeping a single cancellation handle per selection; there
is no automated test for that race, because it needs a controllable slow
response.  It is exercised indirectly by the end-to-end tests that select files
quickly.
\end{boundary}
